% Generated draft: requires mathematical and source review.
\documentclass[11pt]{article}
\usepackage[margin=25mm]{geometry}
\usepackage{fontspec}
\setmainfont{Latin Modern Roman}
\usepackage{amsmath,amssymb,mathtools,unicode-math}
\setmathfont{Latin Modern Math}
\usepackage{enumitem,needspace,xurl,hyperref}
\setlength{\parindent}{0pt}
\setlength{\parskip}{5pt}
\setlength{\emergencystretch}{3em}
\newcommand{\N}{\mathbb{N}}
\newcommand{\Z}{\mathbb{Z}}
\newcommand{\Q}{\mathbb{Q}}
\newcommand{\R}{\mathbb{R}}
\newcommand{\C}{\mathbb{C}}
\newcommand{\F}{\mathbb{F}}
\newcommand{\A}{\mathbb{A}}
\newcommand{\PP}{\mathbb{P}}
\newcommand{\E}{\mathbb{E}}
\newcommand{\eps}{\varepsilon}
\newcommand{\dd}{\,\mathrm{d}}
\begin{document}
\section*{UnsolvedMath ID: 30001459}
\textbf{Source problem code:} catalogue entry
\subsection{Definitions and mathematical statement}
A stationary Poisson cylinder process is a Poisson process of cylinders in a Euclidean space whose law is invariant under translations. The supplied record does not specify the ambient dimension, the dimension of the cylinder axes, or the distribution of cylinder cross-sections.

For standard spherical cosine-transform notation, let \(d\geq 2\), let \(S^{d-1}=\{u\in\R^d:\|u\|=1\}\), and let \(\mu\) be a finite Borel measure on \(S^{d-1}\). Its spherical cosine transform is
\[
(\mathcal C\mu)(u)=\int_{S^{d-1}}|\langle u,v\rangle|\,\mu(\dd v),\qquad u\in S^{d-1}.
\]
If the cylinder directions are unoriented lines, their directional distribution can be represented by an even probability measure \(\mu\), meaning \(\mu(A)=\mu(-A)\) for every Borel set \(A\subseteq S^{d-1}\). This representation is conditional on the intended axis model; that model is not specified in the record. The notation above does not impose a dimension or a range of dimensions on the source problem.

\textbf{Source question.} Given finitely sampled intersection intensities \(g\) of a stationary Poisson cylinder process, can its directional distribution be estimated stably through a continuous inversion of the spherical cosine transform?

In the conditional spherical model above, finite directional samples would have the form
\[
\bigl(u_i,g(u_i)\bigr)_{i=1}^{m},\qquad m\in\N,\quad u_i\in S^{d-1}.
\]
The requested task is to reconstruct an estimate of the directional distribution from such finite observations using cosine-transform inversion. The source does not give the equation relating \(g\) to \(\mathcal C\mu\), including any process-intensity or cross-section factors. It also does not specify sampling assumptions, an observation-error model, admissible directional distributions, topologies for continuity, or a stability or convergence criterion. Consequently, this is an unresolved formulation of an estimation problem, not a fully specified existence assertion for a continuous inverse or estimator. In particular, whether ``continuous inversion'' concerns complete transform data or a reconstruction procedure acting on finite samples remains unresolved.
\subsection{Short English statement}
The problem asks whether finitely many intersection-intensity observations of a stationary Poisson cylinder process allow stable estimation of its directional distribution through spherical cosine-transform inversion. The geometric model, observation equation, sampling assumptions, and meaning of stability require clarification. The supplied record labels the problem open; this draft does not independently verify that status.
\subsection{Sources}
\begin{itemize}
\item Local UnsolvedMath registry, supplied complete record: permanent ID \(30001459\), title \emph{Stable Reconstruction of Poisson Cylinder Directions}, exact source problem code \texttt{catalogue entry}. The record belongs to the Oberwolfach Reports Open Problems collection and gives the proposed year as 2010. No individual proposer is supplied.
\item Supplied report citation: \emph{Mathematics and Algorithms in Tomography} (2010), Oberwolfach Reports, pp.\ 1017--1099, DOI \texttt{10.4171/owr/2010/18}. Supplied source URL and literature-source link, labelled ``original OWR report'': \url{https://doi.org/10.4171/owr/2010/18}. No precise problem page or excerpt is supplied.
\item Literature assessment dated exactly \texttt{2026-08-21}: the supplied curation labels the problem open and states that known cosine-transform inversion results and special numerical methods do not resolve the general estimation problem described here. This is supplied dated curation, not an independently verified literature assessment; the background explicitly requests independent verification of time-sensitive status and cited claims.
\item The supplied statement-status label is \texttt{corrected\_verified}. Its accompanying verification note reports only that the intensity symbol was placed in a mathematical delimiter without changing the question. This metadata does not establish mathematical correctness or independent source verification by this draft.
\end{itemize}
\end{document}
